
\(K\)-algèbres filtrées \(K[X_i]_{i\geq1}\longrightarrow A\) en envoyant \(X_1\) sur \(0\) et \(X_i\) sur le \((i-1)\)-ième coefficient de \(\chi\) pour \(i>1\)~; utilisant cet homomorphisme, on voit qu'il suffit de vérifier que le morphisme de foncteurs considéré s'annule pour l'élément \(1+X_2t+\ldots+X_{n+1}t^n+\ldots\). On utilise alors l'endomorphisme \(f\) de \(K[X_i]_{i\geq1}\) défini par \(f(X_i)=X_{i+1}\) pour tout \(i\), et le diagramme commutatif~:
\[
\begin{tikzcd}
\Ghat(K[(X_i)]) \arrow[r] \arrow[d,"\Ghat(f)"'] &
H\circ\operatorname{Gr}(K[(X_i)]) \arrow[d,"H(\operatorname{Gr}(f))"]\\
\Ghat(K[(X_i)]) \arrow[r] &
H\circ\operatorname{Gr}(K[(X_i)]);
\end{tikzcd}
\]
or \(\operatorname{Gr}(f)=\operatorname{Gr}(f_1)\), \(f_1\) étant l'homomorphisme d'augmentation de \(K[X_i]_{i\geq1}\)~; comme ce dernier annule \(\varphi\), il résulte du diagramme analogue écrit pour \(f_1\) que \(H(\operatorname{Gr}(f))\) annule l'image de \(\varphi\) dans \(H\circ\operatorname{Gr}(K[X_i]_{i\geq1})\), d'où le résultat.

Par passage au quotient, on associe donc à tout morphisme de foncteurs~: \(\Ghat\longrightarrow H\circ\operatorname{Gr}\), un morphisme de foncteurs~: \(1+\widehat{\operatorname{Gr}\,\mathord{.}\,}^{+}\longrightarrow H\circ\operatorname{Gr}\)~; par composition avec \(j\), et en utilisant l'isomorphisme \(\varepsilon\), on obtient un morphisme~: \(1+\widehat{\,\mathord{.}\,}^{+}\longrightarrow H\), ce qui définit l'homomorphisme \(\rho'\) cherché. Il est immédiat de vérifier que \(\rho\) et \(\rho'\) sont inverses l'un de l'autre.

\textbf{Théorème 1.13.} \uline{Soit} \(H\) \uline{un foncteur sur} \(\mathcal G\), \uline{à valeurs dans la catégorie des groupes abéliens filtrés séparés complets}, \uline{et commutant aux invariants sous un groupe fini}. \uline{Alors l'application de} \(\Hom_{\underline{\mathrm{Ab}}\,\underline{\mathrm{fil}}}(1+\widehat{\,\mathord{.}\,}^{+},H)\) \uline{dans} \(H(K[X])\) \uline{qui}, \uline{à un morphisme de foncteurs}, \uline{associe sa valeur sur l'élément} 
